\section{问题四，数据采集建议}

基于问题一至问题三的实际建模经验，本文识别出四项关键数据缺口，并基于量化证据论证每条建议对决策质量的提升路径。

\subsection{建议一，利用已有扫码时间戳优化需求方差估计}

\textbf{来源}：问题二的鲁棒性分析发现，需求方差 $\sigma$ 是利润最敏感的单一参数，$\sigma$ 减半则利润 $+135\%$，$\sigma$ 加倍则亏损。

\textbf{问题}：当前仅以历史日销量的残差标准差作为 $\sigma$ 的粗略估计。附件 2 的扫码时间戳（精确到毫秒）已被记录但未被充分利用。

\textbf{建议}：按小时分段聚合销量，计算日内销售节奏曲线，利用日内分时数据的级数（intraday granularity）获得对日总销量方差的更精确估计。同时可在日间动态调整补货策略，若上午销量显著高于预测，下午启动补充订货。

\textbf{数据成本}：零，此建议是对已有数据的深度利用，无需新增采集。

\subsection{建议二，单品级价格弹性实验}

\textbf{来源}：问题三的 $30/30$ 单品加成率触顶上界。

\textbf{问题}：当前 Q3 使用品类级弹性估计 $\beta_c$ 直接作为单品级弹性的近似。品类内不同单品的价格敏感度可能有巨大差异，明星单品（消费者特别青睐的品种）比品类内普通单品的价格弹性更低。

\textbf{建议}：对核心单品（Top $20\%$ 销量，即贡献 $83.9\%$ 销量的品种）在 $2$ 至 $4$ 周内进行小额价格调整实验（涨跌 $10\%$ 至 $15\%$），记录对应销量响应，估计单品级弹性。

\textbf{帮助}：单品级弹性将直接解决 Q3 的边界饱和问题，不同单品获得差异化的最优加成率，总利润有望提升。

\textbf{数据成本}：低。仅需在正常销售过程中改变标签价格，通过 POS 记录销量响应。

\subsection{建议三，气象与节假日标定}

\textbf{来源}：问题一的趋势分解（STL）仅为数据驱动，无法预先获取未来天气信息。

\textbf{问题}：极端天气（暴雨、持续高温）和特定节假日（春节、中秋）的销量异常无法被模型解释或预测。

\textbf{建议}：接入当地气象数据（日最高/最低温度、降水量），并将农历节假日提前标注。将这些信息作为 Prophet/ETS 预测模型的外生回归变量。

\textbf{帮助}：根据问题二的敏感性分析，预测误差每减少 $10\%$，利润增长约 $13.5\%$（线性外推自 $\sigma \times 0.5$ 场景）。

\textbf{数据成本}：低。中国气象局天气数据对公众免费；节假日日期为公开信息。

\subsection{建议四，陈列空间与竞争价格数据}

\textbf{来源}：问题二和三的"销售空间有限"约束均以历史销量上限作为代理。

\textbf{问题}：真实空间约束取决于陈列面积（$\text{m}^2$）、商品体积和货架空置容忍度，而非销量的统计特性。

\textbf{建议}：一是记录各品类的每日陈列面积与空置时间，建立物理空间到最大补货量的映射。二是每周采集周边 $2$ 至 $3$ 家竞争超市的同品类平均零售价，引入竞争约束。

\textbf{帮助}：空间数据可使 Q2-Q3 的约束从"统计代理"升级为"物理约束"，提高可行性。竞争价格数据可建立交叉价格弹性模型，防止垄断定价假设下的过度乐观。

\textbf{数据成本}：低至中。空间记录每次约 5 分钟，竞争价格采集每周约 30 分钟。

\subsection{优先级排序}

\begin{table}[H]
\centering
\caption{数据采集建议优先级}
\label{tab:q4_priority}
\begin{tabular}{lccc}
\hline
优先级 & 建议 & 对应缺陷 & 预期收益 \\
\hline
1 & 日内销售节奏利用 & Q2 $\sigma$ 敏感性 & 高（已有数据） \\
2 & 单品价格弹性实验 & Q3 边界饱和 & 极高 \\
3 & 气象与节假日标定 & Q1/Q2 预测精度 & 中-高 \\
4 & 陈列空间与竞争价格 & Q2/Q3 约束真实性 & 中 \\
\hline
\end{tabular}
\end{table}

\section{模型评价与推广}

\subsection{模型优点}

\begin{itemize}
\item 系统性去趋势，消除了文献中普遍存在的共同时间趋势伪相关问题
\item 预测不确定性内化，报童模型将需求方差直接转换为风险调整决策，优于确定性的点预测-优化分离范式
\item 内生定价，加成率从固定外生参数变为优化变量，发现了远高于传统 $20\%$ 的利润空间
\end{itemize}

\subsection{模型局限性}

\begin{enumerate}
\item \textbf{单品聚类分离度低}：silhouette = $0.17$，单品销售形态差异不显著，聚类结构弱
\item \textbf{花菜类候选稀缺}：仅 $2$ 个可售单品，是品类覆盖约束的结构性瓶颈
\item \textbf{加成率边界饱和}：$30/30$ 单品（Q3）和 $3/6$ 品类（Q2）的加成率触顶上界 $150\%$，被边界截断而非真正收敛
\item 弹性估计基于观测数据，无外生价格变异，存在识别问题
\item 需求正态性假设和品类独立性假设是简化处理
\end{enumerate}

\subsection{模型推广}

本文的报童 $+$ 弹性内生框架适用于任何保鲜期短、残值低、按日补货的易逝品零售场景，例如鲜奶、烘焙品、鲜切水果和鲜花。核心结构（价格弹性估计 → 需求分布预测 → 随机优化 → 鲁棒性检验）可直接迁移，仅需将蔬菜损耗率替换为目标商品的折旧率。

\newpage

\section{参考文献}

\begin{enumerate}
\item 陈妙霞, 朱映辉, 张赞, 等. 基于 SARIMA 模型的生鲜类商品自动定价与补货决策，以蔬菜类商品为例[J]. 数字技术与应用, 2024(10): 147-150.
\item 高佳楠, 台明浪, 崔思雨, 等. 蔬菜类商品动态定价与补货决策研究[J]. 数学建模及其应用, 2024, 13(2): 47-56.
\item 苏茜, 何暄暄, 卢玥, 等. 基于优化模型的蔬菜类商品定价与补货决策[J]. 软件导刊, 2025, 24(6): 87-94.
\item 聂森, 潘萱颖, 曲浩栋. 基于价格弹性的蔬菜类商品自动定价与补货决策[J]. 数学建模及其应用, 2024, 13(2): 66-72.
\item 曹宇轩, 黄瑞, 秦一天. 基于历史数据的蔬菜类商品定价与补货决策模型[J]. 数学建模及其应用, 2024, 13(2): 37-46.
\item 吴萌, 张驰, 王志勇. "蔬菜类商品的自动定价与补货决策"问题解析[J]. 数学建模及其应用, 2024, 13(2): 27-36.
\item 聂宇旋. 基于 ARIMA 预测优化模型的生鲜类商品自动定价与补货策略研究[J]. 商展经济, 2024(5): 19-22.
\item Liao J, Chen J, Zheng W, et al. Clustering Model-based Analysis of Sales in Vegetable Categories for Supermarkets[C]. 2024.
\item Zhao X, Wang H, Wei R, et al. Automated pricing and replenishment decision model based on single-objective optimization[C]. DEAI 2024.
\end{enumerate}
